% The row classes: every table of the temporal database as a Python object.
\begin{tikzpicture}[font=\scriptsize]
\begin{class}[text width=4.2cm]{DictSQLSerializer}{0,0}
\attribute{D : dict}
\operation{serialize()}
\operation{deserialize()}
\end{class}
\begin{class}[text width=4.2cm]{SQLDatabaseInterface}{0,-2.6}
\inherit{DictSQLSerializer}
\attribute{table : str}
\attribute{ident : str}
\operation{select()}
\operation{insert()}
\operation{update()}
\operation{update_all()}
\operation{delete()}
\operation{is_in_db()}
\operation{get_select_statement()}
\operation{get_insert_statement()}
\end{class}
\begin{class}[text width=4.2cm]{Row}{0,-7.0}
\inherit{SQLDatabaseInterface}
\attribute{columns : Column[*]}
\operation{get_id()}
\operation{set_id()}
\end{class}
\begin{class}[text width=3.6cm]{Column}{6.2,-7.0}
\attribute{name : str}
\attribute{alias : str}
\attribute{convert : callable}
\attribute{writable : bool}
\attribute{default}
\end{class}
\composition{Row}{columns}{*}{Column}
\begin{class}[text width=3.8cm]{Kind}{6.2,-1.6}
\attribute{name : str}
\attribute{aliases : tuple}
\attribute{is_stds : bool}
\attribute{map_kind : str}
\attribute{stds_kind : str}
\attribute{table_prefix : str}
\attribute{has_layer : bool}
\operation{table()}
\end{class}
\begin{class}[text width=2.4cm]{BaseRow}{-6.9,-10.6}
\inherit{Row}
\attribute{kind : Kind}
\attribute{name}
\attribute{mapset}
\attribute{layer}
\attribute{creator}
\attribute{ctime}
\attribute{ttype}
\attribute{semantic_type}
\attribute{mtime}
\end{class}
\begin{class}[text width=2.4cm]{RegisterRow}{-4.2,-10.6}
\inherit{Row}
\attribute{kind : Kind}
\attribute{registered_stds}
\end{class}
\begin{class}[text width=3.0cm]{TemporalExtent}{-1.0,-10.6}
\inherit{Row}
\attribute{start_time}
\attribute{end_time}
\operation{temporal_relation()}
\operation{intersect()}
\operation{union()}
\operation{disjoint_union()}
\operation{starts()}
\operation{during()}
\operation{follows()}
\operation{overlaps()}
\end{class}
\begin{class}[text width=3.0cm]{SpatialExtent}{2.6,-10.6}
\inherit{Row}
\attribute{north}
\attribute{south}
\attribute{east}
\attribute{west}
\attribute{top}
\attribute{bottom}
\operation{spatial_relation()}
\operation{spatial_relation_2d()}
\operation{intersect()}
\operation{union()}
\operation{get_area()}
\end{class}
\begin{class}[text width=2.4cm]{MetadataRow}{6.0,-10.6}
\inherit{Row}
\attribute{kind : Kind}
\operation{print_info()}
\end{class}
\begin{class}[text width=3.5cm]{AbsoluteTemporalExtent}{-3.2,-15.4}
\inherit{TemporalExtent}
\end{class}
\begin{class}[text width=3.5cm]{RelativeTemporalExtent}{0.7,-15.4}
\inherit{TemporalExtent}
\attribute{unit}
\end{class}
\begin{class}[text width=3.5cm]{AbsoluteTimeRow}{-3.2,-17.4}
\inherit{AbsoluteTemporalExtent}
\attribute{kind : Kind}
\attribute{granularity}
\attribute{map_time}
\end{class}
\begin{class}[text width=3.5cm]{RelativeTimeRow}{0.7,-17.4}
\inherit{RelativeTemporalExtent}
\attribute{kind : Kind}
\attribute{granularity}
\attribute{map_time}
\end{class}
\begin{class}[text width=2.6cm]{SpatialExtentRow}{4.6,-15.4}
\inherit{SpatialExtent}
\attribute{kind : Kind}
\end{class}
\end{tikzpicture}
